\section{Multiplicity differences in intervals of triple-logarithmic length}
\label{sec:short}
For distinct primitive Dirichlet characters $\chi_1,\ldots,\chi_s$ of conductors
$q_1,\ldots,q_s$, let $m_j(t)$ be the order of $L(s',\chi_j)$ at $s'=1/2+it$.
The conductor-one character gives the Riemann zeta function. Fix $s\ge2$.
For $c\in\mathbb R^s$ put $z_c(t)=\sum_jc_jm_j(t)$ and
$z_{c,\pm}=\max(\pm z_c,0)$. Write $\log_3 T=\log\log\log T$.
Sums below count each ordinate once.

\begin{theorem}\label{th:short}
Fix $\epsilon>0$ and the characters. For all sufficiently large $T$, every
\[
 (2/\pi+\epsilon)\log_3 T\le H\le T
\]
satisfies, simultaneously for both signs and every $c$ with
$\sum_jc_j=0$, $\|c\|_2=1$,
\[
 \sum_{T<t\le T+H}z_{c,\pm}(t)\ \gg_{\epsilon,\chi_1,\ldots,\chi_s}\ H\log\log H .
\]
\end{theorem}
For two characters take $c=(1,-1)/\sqrt2$. Each interval contains both an
ordinate where the first multiplicity is greater and one where the second is
greater. The constant $2/\pi$ is sufficient as a limiting coefficient in
this construction. Its optimality is open here.

\paragraph{Comparison.}
Srinivas proved that for distinct functions in the Selberg class of
equal degree, each direction of excess occurs within distance
$O(\log\log T)$ of every large $T$, using a kernel of double-exponential
decay. Our use of real critical-line ordinates gives a smaller interval for
fixed Dirichlet families and a lower bound for the mass in it.
Banks gives bounds uniform in the character modulus.
Fujii and Bombieri--Perelli obtain much larger total multiplicity counts in
global ranges. Banks and Loftus detect genuinely unshared zeros in longer
windows. In particular, if the multiplicities at a common zero are 2 and
1, our sum counts one unit there. The conjectural scarcity of common zeros
remains open. See \href{https://doi.org/10.4064/aa103-3-1}{Srinivas},
\href{https://arxiv.org/abs/2302.07073}{Banks},
\href{https://matwbn.icm.edu.pl/ksiazki/aa/aa28/aa2847.pdf}{Fujii},
\href{https://doi.org/10.4064/aa-83-3-271-281}{Bombieri--Perelli}, and
\href{https://arxiv.org/abs/2607.22930}{Banks--Loftus}.

\subsection{A positive weight and its prime frequencies}
Let $Q$ be the least common multiple of the conductors and
$C(n)=\sum_jc_j\chi_j(n)$. The induced characters of
$(\mathbb Z/Q\mathbb Z)^\times$ are distinct by uniqueness of the primitive
inducer. Character orthogonality gives the mean square of $C(v)$ over the
group as $1$. Choose $v$ with $|C(v)|\ge1$, and choose $\theta$ with
$e^{i\theta}C(v)=|C(v)|$. Put $Y=\lambda\log H$ for a sufficiently
small fixed $\lambda>0$. Take all primes $p\le Y$ congruent to $v$ modulo
$Q$ and write
\[
 w_p=\frac{\log p}{\sqrt p},\qquad V^2=\sum w_p^2
       \sim \frac{(\log Y)^2}{2\varphi(Q)}.
\]
The asymptotic follows by partial summation from the prime number theorem
in the finitely many fixed progressions. Form the real trigonometric sums
\[
 A(t)=\Im\prod_p(1+i(w_p/V)\cos(t\log p+\theta)),\qquad
 P_\pm(t)=1\pm e^{-1/2} A(t).
\]
On the real axis $|A|\le\prod_p\exp(w_p^2/(2V^2))=\sqrt e$, so
$0\le P_\pm\le2$. Extend their finite exponential expansions to complex
$t$. Write any one of these two as $P(t)=\sum_R b_R e^{it\log R}$,
where the positive rationals $R$ are in lowest terms.
Unique factorization gives $b_1=1$ and
$b_p=\pm e^{-1/2}e^{i\theta}w_p/(2V)$ at the chosen primes.
Every other nonconstant frequency, apart from their reciprocals, has at
least two different primes with nonzero exponents and cannot be a prime
power or its reciprocal. Also $b_{1/R}=\overline{b_R}$. Both the numerator
and denominator of $R$ are at most $D=\prod_{p\le Y}p\le e^{2Y}$ for
large $Y$. Expanding the product and using Cauchy--Schwarz gives
$\sum_R|b_R|\le 1+\prod_p(1+w_p/V)\le1+\exp(\sqrt{\pi(Y)})$.

\subsection{An explicit-formula test localized in height}
Choose $0<a<\pi$ and $0<\kappa<1/2$ such that
$a(1/2-\kappa)(2/\pi+\epsilon)>1$. Set $B=\kappa H$, $X=T+H/2$.
Let $k(t)$ be $e^{-\cosh(at)}$ divided by its real integral, and put
\[
 K_B(t)=\frac{1}{2B}\int_{-B}^{B}k(t-y)\,dy,\qquad
 h(t)=P(t)K_B(t-X).
\]
Thus $0\le h(t)\le1/B$ on $\mathbb R$. Use the Fourier convention
$\widehat h(u)=\int_{\mathbb R}h(t)e^{-itu}\,dt$.
For fixed $b$ with $1/2<b<\pi/(2a)$, shifting the contour to $\Im t=\pm b$
gives $|\widehat k(u)|\ll_{a,b}e^{-b|u|}$ since
$|e^{-\cosh(a(t+ib))}|=e^{-\cos(ab)\cosh(at)}$.
Consequently
\[
 \widehat h(u)=\sum_R b_R e^{-iX(u-\log R)}
       \widehat k(u-\log R)\operatorname{sinc}(B(u-\log R)),
 \qquad \operatorname{sinc}(w)=\sin(w)/w,
\]
with value $1$ at zero. Resonant frequencies $u=\log R$ contribute exactly
$b_R$ independently of the center $X$.

For a primitive $\chi_j$ write $a_j=(1-\chi_j(-1))/2$ and
$\psi=\Gamma'/\Gamma$. Under Dirichlet GRH the standard Guinand--Weil
formula in this normalization reads
\begin{align*}
 2\pi\sum_t m_j(t)h(t)
   ={}& 2\pi\mathbf1_{q_j=1}\{h(i/2)+h(-i/2)\}\\
 &+\int_{\mathbb R}h(t)\left[\log(q_j/\pi)+
       \Re\psi\big((1/2+a_j+it)/2\big)\right]dt\\
 &-2\Re\sum_{n\ge2}\frac{\Lambda(n)\chi_j(n)}{\sqrt n}\widehat h(\log n).
\end{align*}
The real-valued test need not be even. The formula follows by integrating the
completed logarithmic derivative on the two sides of the critical strip,
using its functional equation on the left. The strip of decay just proved
justifies both contour shifts and absolute summation on the prime side.
For the explicit-formula convention see
\href{https://doi.org/10.1093/qmath/haab041}{\emph{Fourier optimization and quadratic forms}, section 5}.

After multiplication by $c_j$ and summation, the archimedean contribution
is $O(1)$ uniformly. Indeed, the common $\log(|t|/2)$ from Stirling's
formula cancels since $\sum c_j=0$, the combined density is bounded near
zero, and $\int |h|\le2$. The pole contribution is $o(1)$: on
$\Im t=\pm1/2$, $|P(t)|\le D^{1/2}\sum|b_R|$, whereas each translated kernel
in $h(\pm i/2)$ has real displacement at least $X-B\ge T$.

On the prime side the total of the off-diagonal terms is
$O(e^{C Y}/B)$ for a fixed $C$. Here are bounds also covering rational
frequencies close to an integer. For $n\le 2D$ and $n\ne R$, the numerator
of $n-R$ in the common denominator is a nonzero integer, hence
$|\log n-\log R|\ge 1/(2D^2)$. The sinc factor costs at most $2D^2/B$,
and $\sum_{n\le2D}(\log n)/\sqrt n\ll D^{1/2}\log(2D)$.
For $n>2D$ the logarithmic difference exceeds $\log2$ and the Fourier
bound gives $O(B^{-1}D^b n^{-b})$; the resulting sum converges since
$b>1/2$. Summing these bounds with $|b_R|$ proves the assertion, uniformly
in the phases and the vector $c$. Choose $\lambda$ so that $C\lambda<1/2$.
The prime diagonal now gives
\[
 2\pi\sum_t z_c(t)P_\pm(t)K_B(t-X)
    =\mp e^{-1/2}|C(v)|V+O(1).
\]
In particular the two positive test functions have signed integrals of
opposite signs and size $\gg\log\log H$.

\subsection{Removing the tails}
The number of zeros of any fixed primitive function in a unit interval
at height $u$, with multiplicity, is $O(\log(2+|u|))$ by the
Riemann--von Mangoldt formula (see also
\href{https://doi.org/10.1090/mcom/3599}{\emph{Counting zeros of Dirichlet
$L$-functions}}). Put $L=H/2-B$. In the $j$th unit interval outside either
end of $(T,T+H]$, the maximal kernel value is
$O(\exp(-c_a e^{a(L+j)}))$, $j\ge0$. This bound also covers a zero
exactly at the excluded endpoint $T$. Summing the unit-interval bounds,
using $H\le T$, gives a tail of size
\[
 O\!\left(\log(2+T)\exp(-c'_a e^{aL})\right)=o(1).
\]
Indeed $e^{aL}\ge(\log\log T)^{1+\eta}$ for some $\eta>0$.
The same estimates apply to negative ordinates. Since $0\le h\le1/B$,
each sign inside the interval has mass $\gg B V$. This proves
Theorem~\ref{th:short}.\hfill$\square$

\paragraph{A single-prime illustration.}
For $\zeta$ and the Dirichlet beta function $L(s,\chi_{-4})$, use the
unnormalized difference of their multiplicities. The unamplified conductor
term is $-\log 4$. Using $P=1\pm\cos(t\log7)$, the two limiting
signed readings after multiplication by $2\pi$ are
$-\log4\mp 2\log7/\sqrt7$. The absolute prime term exceeds $\log4$,
so this one prime already certifies both signs for sufficiently large
windows of the stated length. The growing family of primes in the proof
supplies the factor $\log\log H$.
